\section*{Chapter \ref{ch:b2:ligand-field} --- Ligand Field Theory and Colour}

\begin{solution}{exo:b2:ligand-field:1}
$d^3$: $t_{2g}^3$, CFSE $-1.2\Delta_o$. $d^8$: $t_{2g}^6e_g^2$, CFSE $-2.4 + 1.2 = -1.2\Delta_o$.
\end{solution}

\begin{solution}{exo:b2:ligand-field:2}
\ce{[Fe(H2O)6]^2+}, $d^6$ \omterm{def:b2:ligand-field:spin}{high spin} $t_{2g}^4e_g^2$: four. \ce{[Fe(CN)6]^4-}, \omterm{def:b2:ligand-field:spin}{low spin}
$t_{2g}^6$: none (\omterm{def:b2:diatomic-mos:magnetism}{diamagnetic}).
\end{solution}

\begin{solution}{exo:b2:ligand-field:3}
600 nm is orange light; the complex looks blue (greenish blue).
\end{solution}

\begin{solution}{exo:b2:ligand-field:4}
$\ce{Cl-} < \ce{H2O} < \ce{NH3} < \ce{CN-}$.
\end{solution}

\begin{solution}{exo:b2:ligand-field:5}
Water: $\Delta_o < P$, \omterm{def:b2:ligand-field:spin}{high spin}, $t_{2g}^3e_g^2$, five unpaired electrons. Cyanide: $\Delta_o > P$,
\omterm{def:b2:ligand-field:spin}{low spin}, $t_{2g}^5$, one unpaired electron.
\end{solution}

\begin{solution}{exo:b2:ligand-field:6}
In the tetrahedral \ce{[CoCl4]^2-} the splitting is about $4/9$ of an octahedral one, and
chloride is a weaker ligand than water: the d--d band moves to longer wavelengths (orange
to red), and the complex looks blue. Tetrahedral complexes also absorb more strongly,
hence the deep colour.
\end{solution}

\begin{solution}{exo:b2:ligand-field:7}
\ce{Ni^2+} is $d^8$. Square planar: four lower orbitals hold the eight electrons in pairs,
$\mathrm d_{x^2-y^2}$ stays empty: \omterm{def:b2:diatomic-mos:magnetism}{diamagnetic}. Tetrahedral: $e^4t_2^4$, two unpaired electrons
in the $t_2$ set.
\end{solution}

\begin{solution}{exo:b2:ligand-field:8}
$\tilde\nu = 1/(500 \times 10^{-7}\,\text{cm}) = \qty{20000}{cm^{-1}}$; $\Delta_o =
0.1196/(500 \times 10^{-9}) = \qty{2.39e5}{J/mol}$, \qty{239}{kJ/mol}.
\end{solution}

\begin{solution}{exo:b2:ligand-field:9}
\ce{Zn^2+} is $d^{10}$ and \ce{Sc^3+} $d^0$: no \omterm{def:b2:ligand-field:d-d}{d--d transition} is possible (no empty d level
for \ce{Zn^2+}, no d electron for \ce{Sc^3+}), so no absorption in the visible.
\end{solution}

\begin{solution}{exo:b2:ligand-field:10}
Co(III) $d^6$ + six ammonia lone pairs: 18 electrons. Twelve fill the six \omterm{def:b2:diatomic-mos:bonding}{bonding orbitals},
six fill the nonbonding $t_{2g}$ (ammonia gives a field strong enough for \omterm{def:b2:ligand-field:spin}{low spin} with
cobalt(III)); $e_g^*$ is empty. All electrons are paired: \omterm{def:b2:diatomic-mos:magnetism}{diamagnetic}.
\end{solution}

\begin{solution}{exo:b2:ligand-field:11}
\ce{CO} is a good $\sigma$ donor and, above all, a $\pi$ acceptor: its empty $\pi^*$ orbitals
lower the $t_{2g}$ level by \omterm{def:b2:ligand-field:pi-ligands}{back-donation}, widening $\Delta_o$. \ce{F-} is a $\pi$ donor, which
raises $t_{2g}$ and narrows $\Delta_o$. The point-charge model sees only the charge and gets the
order wrong.
\end{solution}

\begin{solution}{exo:b2:ligand-field:12}
Without a field, the hydration enthalpies would follow a smooth line (ions shrinking along
the series). Water gives high-spin aqua complexes whose CFSE is zero for $d^0$, $d^5$, $d^{10}$
and largest for $d^3$ and $d^8$: the extra stabilisation adds two humps, peaking near
\ce{V^2+} and \ce{Ni^2+}, above the line through \ce{Ca^2+}, \ce{Mn^2+}, \ce{Zn^2+}.
\end{solution}

\begin{solution}{pb:b2:ligand-field:1}
\textbf{1.} $d^3$.
\textbf{2.} $t_{2g}^3$, one electron in each of the three orbitals, parallel spins.
\textbf{3.} No: three electrons fit in $t_{2g}$ without pairing.
\textbf{4.} $-1.2\Delta_o$.
\textbf{5.} Three.
\textbf{6.} The CFSE is large and the $e_g^*$ orbitals, which point at the ligands, are
empty: removing or adding a ligand costs much of that stabilisation.
\textbf{7.} Six oxide ions at the corners of a slightly distorted octahedron.
\textbf{8.} $1/(556 \times 10^{-7}) = \qty{17990}{cm^{-1}}$, about \qty{1.80e4}{cm^{-1}}.
\textbf{9.} $0.1196/(556 \times 10^{-9}) = \qty{2.15e5}{J/mol}$: \qty{215}{kJ/mol}.
\textbf{10.} Yellow-green (556 nm) and violet (405 nm).
\textbf{11.} Red, with a little blue: the deep red of ruby.
\textbf{12.} A $d^3$ ion has more than one excited configuration with an electron in $e_g$;
electron repulsion gives them different energies, hence several bands (treated in the
Year 3 volume).
\textbf{13.} $0.1196/(620 \times 10^{-9}) = \qty{1.93e5}{J/mol}$, \qty{193}{kJ/mol}
(\qty{16130}{cm^{-1}}).
\textbf{14.} Ruby: its splitting is larger.
\textbf{15.} To orange-red: red is now absorbed, and green (with some blue) is transmitted.
\textbf{16.} In beryl the chromium sites are slightly larger and the oxide ions farther,
and the surroundings differ: a weaker field.
\textbf{17.} No: a $d^3$ ion keeps three unpaired electrons in any octahedral field.
\textbf{18.} $d^9$.
\textbf{19.} The aqua complex absorbs in the red and orange by a \omterm{def:b2:ligand-field:d-d}{d--d transition}: the
transmitted light is blue.
\textbf{20.} It turns white: without water ligands around copper (sulfate takes their
place), the field is weaker and the band moves into the infrared.
\textbf{21.} Water binds to copper again and the blue aqua complex forms.
\textbf{22.} At shorter wavelength: ammonia, above water in the series, gives a larger
splitting (the deep blue-violet of the ammine complex).
\textbf{23.} No: $t_{2g}^6e_g^3$ is the only configuration.
\textbf{24.} Yes: one unpaired electron.
\textbf{25.} $\boldsymbol{\Delta_o \approx \qty{1.8e4}{cm^{-1}}}$, about
$\boldsymbol{\qty{215}{kJ/mol}}$.
\end{solution}
